\documentclass[10pt]{extarticle}
\usepackage[a4paper,margin=1.5cm]{geometry}
\usepackage[T1]{fontenc}
\usepackage{amsmath,amssymb}
\usepackage{graphicx}
\usepackage[export]{adjustbox}
\usepackage{multicol}
\usepackage{enumitem}
\usepackage{booktabs,array,tabularx}
\setlength{\columnsep}{0.7cm}
\setlength{\parindent}{0pt}
\setlength{\parskip}{3pt}
\pagestyle{empty}
\makeatletter
% figure registry (figdefs.tex is generated by the extract step)
\newcommand{\qfigdef}[4]{\@namedef{qf@file@#1}{#2}\@namedef{qf@w@#1}{#3}\@namedef{qf@wide@#1}{#4}}
\newcommand{\qf@check}[1]{\@ifundefined{qf@file@#1}{\PackageError{score2}{Unknown figure id #1}{}}{}}
\newcommand{\qf@img}[4][]{\includegraphics[width=#2\dimexpr\@nameuse{qf@w@#4}pt\relax,max width=#3,max height=0.45\textheight,#1]{\@nameuse{qf@file@#4}}}
% Full-width questions: finalize wraps every question that contains a wide figure in \qfullbegin...\qfullend,
% so the WHOLE question (stem, figures, options) is set across the page in reading order, then the two columns
% resume. The question's head - its number, stem and first figure (\qfullsplit marks the end of the head) -
% is kept on one page, so a number is never left at a page foot with its figure overleaf; the rest may break
% across pages normally, so no big gaps are left. Without a \qfullsplit, the whole question is kept together
% if it is under 45% of a page. Never write these yourself.
\newif\ifqf@full
\newif\ifqf@boxed
\newsavebox{\qf@box}
\newcommand{\qf@open}{\global\qf@fulltrue\par\global\qf@boxedtrue
  \setbox\qf@box\vbox\bgroup\hsize\textwidth\linewidth\textwidth}
\newcommand{\qf@emit}[1]{\par\egroup\global\qf@boxedfalse
  \ifdim\dimexpr\ht\qf@box+\dp\qf@box\relax<#1\textheight\noindent\usebox{\qf@box}\par\else\unvbox\qf@box\fi}
\newcommand{\qfullsplit}{\ifqf@boxed\qf@emit{0.85}\fi}
\newcommand{\qf@close}{\ifqf@boxed\qf@emit{0.45}\fi\global\qf@fullfalse}
\newcommand{\qfullbegin}{\end{multicols}\qf@open}
\newcommand{\qfullend}{\qf@close\begin{multicols}{2}}
\newcommand{\qfullstart}{\qf@open}
\newcommand{\qfulldone}{\qf@close}
% \qfig[scale]{id}: block figure, centred, directly where it is written (size set by the tool).
% Fallback if a wide figure is not inside a full-width question: break out of the columns in place - no floats.
\newcommand{\qfig}[2][1]{\qf@check{#2}%
  \ifqf@full\qf@blk{#1}{#2}\else
  \ifnum\@nameuse{qf@wide@#2}=1 \expandafter\@firstoftwo\else\expandafter\@secondoftwo\fi
  {\end{multicols}\begin{center}\qf@img{#1}{\textwidth}{#2}\end{center}\begin{multicols}{2}}%
  {\qf@blk{#1}{#2}}\fi}
\newcommand{\qf@blk}[2]{\par\smallskip{\centering\qf@img{#1}{\linewidth}{#2}\par}\smallskip}
% \qfigi[scale]{id}: inline figure (inside option lists / tables / minipages); never breaks columns.
\newcommand{\qfigi}[2][1]{\qf@check{#2}\qf@img[valign=c]{#1}{\linewidth}{#2}}
\makeatother
% structure macros - lint reads these to check completeness and figure placement
\newcommand{\papersec}[1]{\par\bigskip{\centering\Large\bfseries PAPER-#1\par}\medskip}
\newcommand{\partsec}[2]{\par\medskip\noindent{\large\bfseries PART-#1: #2}\par\smallskip}
\newcommand{\sect}[1]{\par\smallskip\noindent\textbf{#1}\par}
\newcommand{\qn}[1]{\par\smallskip\noindent\textbf{#1.}~}
\newcommand{\paraq}[2]{\par\smallskip\noindent\textbf{Paragraph for Questions #1 and #2}\par}
\newenvironment{qoptions}{\begin{itemize}[leftmargin=1.8em,itemsep=1pt,topsep=2pt,label={}]}{\end{itemize}}
\newcommand{\opt}[2]{\item[(#1)] #2}
% figure registry (generated by score2_tools.py finalize)
\qfigdef{p1_phy_s1_q02_1}{Score Advanced - Test 31 (04 May 2024)-img/p1_phy_s1_q02_1.png}{108.3}{0}
\qfigdef{p1_phy_s1_q03_1}{Score Advanced - Test 31 (04 May 2024)-img/p1_phy_s1_q03_1.png}{171.7}{0}
\qfigdef{p1_phy_s2_q06_1}{Score Advanced - Test 31 (04 May 2024)-img/p1_phy_s2_q06_1.png}{74.2}{0}
\qfigdef{p1_phy_s2_q07_1}{Score Advanced - Test 31 (04 May 2024)-img/p1_phy_s2_q07_1.png}{92.6}{0}
\qfigdef{p1_phy_s2_q08_1}{Score Advanced - Test 31 (04 May 2024)-img/p1_phy_s2_q08_1.png}{128.3}{0}
\qfigdef{p1_phy_s2_q09_1}{Score Advanced - Test 31 (04 May 2024)-img/p1_phy_s2_q09_1.png}{227.2}{0}
\qfigdef{p1_phy_s3_q01_1}{Score Advanced - Test 31 (04 May 2024)-img/p1_phy_s3_q01_1.png}{92.3}{0}
\qfigdef{p1_phy_s3_q01_2}{Score Advanced - Test 31 (04 May 2024)-img/p1_phy_s3_q01_2.png}{92.5}{0}
\qfigdef{p1_phy_s3_q01_3}{Score Advanced - Test 31 (04 May 2024)-img/p1_phy_s3_q01_3.png}{92.5}{0}
\qfigdef{p1_phy_s3_q03_1}{Score Advanced - Test 31 (04 May 2024)-img/p1_phy_s3_q03_1.png}{134.2}{0}
\qfigdef{p1_phy_s3_q04_1}{Score Advanced - Test 31 (04 May 2024)-img/p1_phy_s3_q04_1.png}{45.1}{0}
\qfigdef{p1_phy_s3_q05_1}{Score Advanced - Test 31 (04 May 2024)-img/p1_phy_s3_q05_1.png}{149.6}{0}
\qfigdef{p1_phy_s4_q01_1}{Score Advanced - Test 31 (04 May 2024)-img/p1_phy_s4_q01_1.png}{93.5}{0}
\qfigdef{p1_chem_s1_q04_D}{Score Advanced - Test 31 (04 May 2024)-img/p1_chem_s1_q04_D.png}{44.9}{0}
\qfigdef{p1_chem_s1_q05_R}{Score Advanced - Test 31 (04 May 2024)-img/p1_chem_s1_q05_R.png}{164.1}{0}
\qfigdef{p1_chem_s1_q05_A}{Score Advanced - Test 31 (04 May 2024)-img/p1_chem_s1_q05_A.png}{46.1}{0}
\qfigdef{p1_chem_s1_q05_B}{Score Advanced - Test 31 (04 May 2024)-img/p1_chem_s1_q05_B.png}{65.1}{0}
\qfigdef{p1_chem_s1_q05_C}{Score Advanced - Test 31 (04 May 2024)-img/p1_chem_s1_q05_C.png}{42.4}{0}
\qfigdef{p1_chem_s1_q05_D}{Score Advanced - Test 31 (04 May 2024)-img/p1_chem_s1_q05_D.png}{76.8}{0}
\qfigdef{p1_chem_s2_q09_1}{Score Advanced - Test 31 (04 May 2024)-img/p1_chem_s2_q09_1.png}{132.5}{0}
\qfigdef{p1_chem_s3_q03_1}{Score Advanced - Test 31 (04 May 2024)-img/p1_chem_s3_q03_1.png}{49.2}{0}
\qfigdef{p1_chem_s3_q04_1}{Score Advanced - Test 31 (04 May 2024)-img/p1_chem_s3_q04_1.png}{216.3}{0}
\qfigdef{p1_chem_s4_q02_A}{Score Advanced - Test 31 (04 May 2024)-img/p1_chem_s4_q02_A.png}{64.7}{0}
\qfigdef{p1_chem_s4_q02_B}{Score Advanced - Test 31 (04 May 2024)-img/p1_chem_s4_q02_B.png}{95.1}{0}
\qfigdef{p2_phy_s1_q01_1}{Score Advanced - Test 31 (04 May 2024)-img/p2_phy_s1_q01_1.png}{126.4}{0}
\qfigdef{p2_phy_s1_q03_1}{Score Advanced - Test 31 (04 May 2024)-img/p2_phy_s1_q03_1.png}{87.0}{0}
\qfigdef{p2_phy_s1_q04_1}{Score Advanced - Test 31 (04 May 2024)-img/p2_phy_s1_q04_1.png}{154.8}{0}
\qfigdef{p2_phy_s1_q05_1}{Score Advanced - Test 31 (04 May 2024)-img/p2_phy_s1_q05_1.png}{86.5}{0}
\qfigdef{p2_phy_s1_q07_1}{Score Advanced - Test 31 (04 May 2024)-img/p2_phy_s1_q07_1.png}{76.8}{0}
\qfigdef{p2_phy_s2_para8_9_1}{Score Advanced - Test 31 (04 May 2024)-img/p2_phy_s2_para8_9_1.png}{121.5}{0}
\qfigdef{p2_phy_s2_para10_11_1}{Score Advanced - Test 31 (04 May 2024)-img/p2_phy_s2_para10_11_1.png}{121.5}{0}
\qfigdef{p2_phy_s2_q11_A}{Score Advanced - Test 31 (04 May 2024)-img/p2_phy_s2_q11_A.png}{87.2}{0}
\qfigdef{p2_phy_s2_q11_B}{Score Advanced - Test 31 (04 May 2024)-img/p2_phy_s2_q11_B.png}{92.6}{0}
\qfigdef{p2_phy_s2_q11_C}{Score Advanced - Test 31 (04 May 2024)-img/p2_phy_s2_q11_C.png}{92.5}{0}
\qfigdef{p2_phy_s2_q11_D}{Score Advanced - Test 31 (04 May 2024)-img/p2_phy_s2_q11_D.png}{92.6}{0}
\qfigdef{p2_phy_s3_q01_1}{Score Advanced - Test 31 (04 May 2024)-img/p2_phy_s3_q01_1.png}{47.8}{0}
\qfigdef{p2_phy_s3_q03_1}{Score Advanced - Test 31 (04 May 2024)-img/p2_phy_s3_q03_1.png}{136.5}{0}
\qfigdef{p2_phy_s3_q06_1}{Score Advanced - Test 31 (04 May 2024)-img/p2_phy_s3_q06_1.png}{102.7}{0}
\qfigdef{p2_chem_s1_q01_1}{Score Advanced - Test 31 (04 May 2024)-img/p2_chem_s1_q01_1.png}{88.7}{0}
\qfigdef{p2_chem_s1_q03_A}{Score Advanced - Test 31 (04 May 2024)-img/p2_chem_s1_q03_A.png}{75.4}{0}
\qfigdef{p2_chem_s1_q03_B}{Score Advanced - Test 31 (04 May 2024)-img/p2_chem_s1_q03_B.png}{115.6}{0}
\qfigdef{p2_chem_s1_q03_C}{Score Advanced - Test 31 (04 May 2024)-img/p2_chem_s1_q03_C.png}{95.5}{0}
\qfigdef{p2_chem_s1_q03_D}{Score Advanced - Test 31 (04 May 2024)-img/p2_chem_s1_q03_D.png}{138.2}{0}
\qfigdef{p2_chem_s1_q04_1}{Score Advanced - Test 31 (04 May 2024)-img/p2_chem_s1_q04_1.png}{119.0}{0}
\qfigdef{p2_chem_s1_q05_A}{Score Advanced - Test 31 (04 May 2024)-img/p2_chem_s1_q05_A.png}{134.2}{0}
\qfigdef{p2_chem_s1_q05_B}{Score Advanced - Test 31 (04 May 2024)-img/p2_chem_s1_q05_B.png}{167.2}{0}
\qfigdef{p2_chem_s1_q05_C}{Score Advanced - Test 31 (04 May 2024)-img/p2_chem_s1_q05_C.png}{128.1}{0}
\qfigdef{p2_chem_s1_q05_D}{Score Advanced - Test 31 (04 May 2024)-img/p2_chem_s1_q05_D.png}{186.3}{0}
\qfigdef{p2_chem_s3_q03_S}{Score Advanced - Test 31 (04 May 2024)-img/p2_chem_s3_q03_S.png}{245.2}{0}
\qfigdef{p2_chem_s3_q04_1}{Score Advanced - Test 31 (04 May 2024)-img/p2_chem_s3_q04_1.png}{133.2}{0}
\raggedcolumns
\makeatletter\@beginparpenalty=10000\makeatother
\newcolumntype{L}[1]{>{\raggedright\arraybackslash}p{#1}}
\begin{document}
\begin{center}{\Large\bfseries Score Advanced --- Test 31 (04 May 2024)}\end{center}
\begin{multicols}{2}
\raggedcolumns
\papersec{1}
\partsec{1}{PHYSICS}
\sect{SECTION-I (i)}
\qn{1} Free $^{238}$U nuclei kept in a train emit alpha particles. When the train is stationary and a uranium nucleus decays, a passenger measures that the separation between the alpha particle and the recoiling nucleus becomes x in time t after the decay. If a decay takes place when the train is moving at a uniform speed v, the distance between the alpha particle and the recoiling nucleus at a time t after the decay, as measured by the passenger will be :-
\begin{qoptions}
\opt{A}{x + vt}
\opt{B}{x $-$ vt}
\opt{C}{x}
\opt{D}{depends on the direction of the train.}
\end{qoptions}
\qn{2} Two identical spherical masses are kept at some distance as shown. Potential energy when a mass m is taken from surface of one sphere to the other.\par
\qfig{p1_phy_s1_q02_1}
\begin{qoptions}
\opt{A}{Increases continuously.}
\opt{B}{Decreases continuously.}
\opt{C}{First increase then decreases.}
\opt{D}{First decreases then increases.}
\end{qoptions}
\qn{3} A uniform ring of mass m and radius a and a sphere of mass M and radius a are released in free space. Consider only gravitational force between sphere and ring. Find acceleration of ring at given separation.\par
\qfig{p1_phy_s1_q03_1}
\begin{qoptions}
\opt{A}{$\dfrac{\sqrt3\,\mathrm{GM}}{8a^3}$}
\opt{B}{$\dfrac{\sqrt3\,\mathrm{GM}}{4a^2}$}
\opt{C}{$\dfrac{\sqrt3\,\mathrm{GM}}{8a^2}$}
\opt{D}{$\dfrac{\mathrm{GM}}{\sqrt3\,a^2}$}
\end{qoptions}
\qn{4} A car drives along a road that has small depressions regularly spaced about 8.0 m apart. When four 80 kg passengers enter the car, it sinks down by 1.6 cm. At approximately what speed the car might travel so that travelling in the vehicle becomes very uncomfortable?
\begin{qoptions}
\opt{A}{5$\pi$ m/s}
\opt{B}{10$\pi$ m/s}
\opt{C}{2$\pi$ m/s}
\opt{D}{20$\pi$ m/s}
\end{qoptions}
\qn{5} Three tuning forks in the audible range produce beats: B and C produce beats at 7 Hz, A and C produce beats at 8 Hz. The frequency of B is 6\% higher than that of A.
\begin{qoptions}
\opt{A}{Frequency of tuning fork A is 265 Hz.}
\opt{B}{Frequency of tuning fork B is 265 Hz.}
\opt{C}{Frequency of tuning fork C is 250 Hz.}
\opt{D}{None of the above}
\end{qoptions}
\sect{SECTION-I (ii)}
\qn{6} Two identical thin homogeneous nichrome wires are mounted into the side non-conducting walls of vertical communicating vessels (vessel cross-sectional area S = 100 cm$^2$ and 3S). The length of the wires is equal to the height of the vessels. An ideal battery with a voltage at the terminals U = 46.2 V is connected to the lower ends of the wires. The vessels are filled upto a level 50 cm and covered with light and thin conductive pistons. The pistons do not allow water to pass through, have contact with the wires and can move without friction inside the vessels without leaving them, thanks to stoppers.\par
An ammeter is connected to the pistons using flexible wires (the length of the wires allows the pistons to move freely). The installation diagram is shown in the figure. When a non-conducting weight weighing 1 kg is placed on a small piston, the ammeter shows a current strength of $I_1$ = 2.31 A. If 3 such weights are placed on a large piston, the ammeter shows a current strength of $I_2$ = 2.1 A.\par
We can assume that the connecting tube has a negligibly small cross-section. The conductivity of water is also very low. Water density $\rho$ = 1000 kg/m$^3$.\par
\qfig{p1_phy_s2_q06_1}
\begin{qoptions}
\opt{A}{The resistance of the ammeter R$_\mathrm a$ is 1 ohm.}
\opt{B}{The reading of ammeter when no weight is placed on the pistons is approx. 2.2 A.}
\opt{C}{The maximum current through the ammeter is approx. 3.22 A.}
\opt{D}{The minimum current through the ammeter is approx. 1.5 A.}
\end{qoptions}
\qn{7} Three tigers simultaneously begin to move along a horizontal surface with constant speeds. The velocity of the first tiger at any moment of time is directed towards the second, the velocity of the second towards the third, and the velocity of the third towards the first. At the initial moment of time, tigers form a right triangle with legs equal to L (Figure). Consider the size of the tigers to be much smaller than L. The speed of the first tiger is v, and the speeds of the second and third tigers v$_2$ and v$_3$ are such that during the movement the angles in the triangle ABC formed by the tigers remain constant.\par
Let's introduce a coordinate system as shown in the figure. The origin of coordinates coincides with the position of the first tiger at the moment of start (point A).\par
\qfig{p1_phy_s2_q07_1}
\begin{qoptions}
\opt{A}{The time after which the tigers will meet is $\dfrac{\mathrm L}{\mathrm v}$.}
\opt{B}{The speed of second tiger v$_2$ is $\dfrac{\mathrm v}2$.}
\opt{C}{The speed of third tiger v$_3$ is $\dfrac{\mathrm v}2$.}
\opt{D}{The coordinates of the point where the tigers will meet is $\left(\dfrac{4\mathrm L}5,\dfrac{\mathrm L}5\right)$}
\end{qoptions}
\qn{8} As shown in figure, a plano-convex lens has a refractive index n and is placed in the air. The radius of the lens aperture is R. Take a point F on the principal optical axis of the lens, OF = f. When parallel light is incident along the principal optical axis, all the light rays converge at point F.\par
\qfig{p1_phy_s2_q08_1}
\begin{qoptions}
\opt{A}{The shape of the convex surface of the lens should be hyperbola.}
\opt{B}{The shape of the convex surface of the lens should be parabola.}
\opt{C}{The maximum value of aperture R of the lens is nf.}
\opt{D}{The maximum value of aperture R of the lens is $\sqrt{n^2-1}\,$f.}
\end{qoptions}
\qn{9} An ice hemisphere with radius R and refractive index n lies on a warm flat table and melts slowly. The rate of heat transfer from the table to the ice is proportional to the area of contact between them. It is known that the ice hemisphere completely melts in time T$_0$. Throughout the process, a laser beam is incident on the ice from above. The beam is vertically incident at a distance of R/2 from the axis of symmetry (see figure).\par
Assume that the temperature of the ice and the surrounding atmosphere are 0$^\circ$C and remains constant during the melting process. The laser beam does not transfer energy to the ice. The melting water immediately flows off the table, and the ice does not move along the table.\par
\qfig{p1_phy_s2_q09_1}
\begin{qoptions}
\opt{A}{The value of x$_0$ (point on the table where the laser beam strikes at t = 0) is $\dfrac{2\mathrm R}{\left(\sqrt{12n^2-3}+1\right)}$.}
\opt{B}{The value of x$_0$ (point on the table where the laser beam strikes at t = 0) is $\dfrac{2\mathrm R}{\left(\sqrt{4n^2-3}+1\right)}$.}
\opt{C}{The height of ice z at any time t is z = R$\left(1-\dfrac{\mathrm t^2}{\mathrm T_0^2}\right)$.}
\opt{D}{The time at which the value of x (point on the table where the laser beam strikes) becomes $\dfrac{\mathrm R}2$ is $\dfrac{\sqrt3}2$T$_0$.}
\end{qoptions}
\sect{SECTION-II}
\qn{1} A metallic rod of mass m and length L can slide without friction on two perpendicular wires. Entire arrangement is located in the horizontal plane. A magnetic field of magnitude B exists perpendicular to this plane in the downward direction. The wire have negligible resistance compared to the rod whose resistance is R. Initially rod is along one of the wire so that one end of it is at the junction of two wires show the rod is given an initial angular speed $\omega$ such that it slide with two ends in contact with the two wires \& just comes to rest in an aligned position with the other wire. The value of $\omega$ is $\dfrac{3\pi\mathrm B^2\mathrm L^2}{\alpha\mathrm{mR}}$. The value of $\alpha$ is :- (Neglect self inductance of system.)\par
\qfig{p1_phy_s3_q01_1}
\qfig{p1_phy_s3_q01_2}
\qfig{p1_phy_s3_q01_3}
\qn{2} Two blocks of metal with heat capacities C$_1$ = 3 J/K and C$_2$ = 6 J/K, are interconnected by a rod of length $\ell$ and cross-sectional area S and fairly low heat conductivity K. The whole system is thermally insulated from the environment. At a moment t = 0 the temperature difference between the two blocks of metal equals ($\Delta$T)$_0$. Assuming the heat capacity of the rod to be negligible, if $\dfrac{\Delta\mathrm T}{\Delta\mathrm T_0}$ = e$^{-\frac{\mathrm{SKt}}{\alpha\ell}}$, then $\alpha$ is\par
\qn{3} A sphere of mass m falls symmetrically on a set of three hemispheres placed in contact with each other on smooth horizontal ground from a height h. Mass of each hemisphere is m/2 and radius of each body is R. Neglect any rebounding of hemisphere. If sphere comes to rest after the simultaneous collision with three hemispheres then ratio of KE of sphere before collision and KE of each hemisphere after collision is p. Find the value of p.\par
\qfig{p1_phy_s3_q03_1}
\qn{4} A syringe is filled with water up to volume 20 cm$^3$. The area of cross-section of the cylinder is 5 cm$^2$. The syringe is held fixed vertically and its 90 gm piston is pushed upward by external agent with constant speed. A water beam coming out of the small nozzle (hole area 1 mm$^2$) has speed 2 m/s. Neglecting friction and viscous nature of water, find the work done by the agent (in 10$^{-2}$ J) in fully emptying the syringe. (Take : g = 10 m/s$^2$)\par
\qfig{p1_phy_s3_q04_1}
\qn{5} A positive charge with a constant surface density $\sigma$ is distributed over the surface of a thin flat non-conducting ring with internal and external radii of r$_1$ and r$_2$, respectively. The center of a non-conducting ball with radius r $<$ r$_1$ coincides with the center of the ring. The sphere is uniformly charged in volume with a positive charge Q. At what values of $\sigma$ (in C/m$^2$) is the resulting field strength vector at points on the surface of the ball located in the plane of the ring is zero? (Take : R = 1 cm, r$_1$ = 2 m, r$_2$ = 3 m and Q = $\pi\times10^{-6}$ C)\par
\qfig{p1_phy_s3_q05_1}
\qn{6} We use potential energy as mgh which is an approximate value, the actual value is mgR$^2\left(\dfrac1{\mathrm R}-\dfrac1{\mathrm R+\mathrm h}\right)$. Where R is radius of earth (6400 km) and h is height from earth's surface. Find maximum height (in km) from earth's surface where the potential energy expression has error of less than 5\% of true value.
\sect{SECTION-IV}
\qn{1} A uniform plank of mass m, free to move in the horizontal direction, is placed on the top of a solid cylinder of mass 2m and radius R. The plank is attached to a fixed wall with the help of a light spring of force constant k. The plank is slightly displaced towards spring and released. There is no slipping anywhere. For any displacement x, amplitude x$_\mathrm o$ and velocity v of the plank, match the column-I with column-II and mark the correct option from the codes given below.\par
\qfig{p1_phy_s4_q01_1}
{\footnotesize\setlength{\tabcolsep}{1.5pt}\noindent\begin{tabular}{|c|L{3.6cm}|c|L{2.3cm}|}
\hline
 & \multicolumn{1}{>{\centering\arraybackslash}p{3.6cm}|}{Column-I} & & \multicolumn{1}{>{\centering\arraybackslash}p{2.3cm}|}{Column-II} \\ \hline
(A) & Velocity of centre of cylinder. & (P) & $\dfrac12$kx$_0^2$ \\ \hline
(B) & Energy of oscillations of the system. & (Q) & $\dfrac{\mathrm v}2$ \\ \hline
(C) & Magnitude of acceleration of the plank. & (R) & $\dfrac78$mv$^2$ + $\dfrac12$kx$^2$ \\ \hline
(D) & Angular frequency of oscillations. & (S) & $\dfrac{4\mathrm{kx}}{7\mathrm m}$ \\ \hline
 & & (T) & $\sqrt{\dfrac{4\mathrm k}{7\mathrm m}}$ \\ \hline
\end{tabular}}\par
\qn{2} Match the following.\par
{\footnotesize\setlength{\tabcolsep}{1.5pt}\noindent\begin{tabular}{|c|L{4.6cm}|c|L{1.8cm}|}
\hline
 & \multicolumn{1}{>{\centering\arraybackslash}p{4.6cm}|}{Column-I (Situation)} & & \multicolumn{1}{>{\centering\arraybackslash}p{1.8cm}|}{Column-II (Values in MKS units)} \\ \hline
(A) & Magnitude of electric field at (1,1,1) if potential is given by V = $-\dfrac1{\mathrm x}-\dfrac1{\mathrm y}-\dfrac1{\mathrm z}$, is & (P) & 0.4 \\ \hline
(B) & The charge flowing through a resistance 6$\Omega$ varies with t as Q = 4t $-$ 2t$^2$. The total heat produced in resistance is & (Q) & $\sqrt3$ \\ \hline
(C) & A milliammeter of range 10mA has a coil of resistance 0.99 $\Omega$. To use it as an ammeter of range 1A, the required shunt must have a resistance of & (R) & 0.01 \\ \hline
(D) & Two unlike point charges of same magnitude 10 mC are placed on x-axis at a separation of 8 cm. If both charges placed symmetrically about origin in external electric field $\vec{\mathrm E}$ = 500 V m$^{-1}$ $\hat{\mathrm i}$, then potential energy of the system may be [Assume potential energy to be zero when dipole is perpendicular to electric field] & (S) & 32 \\ \hline
 & & (T) & None of these \\ \hline
\end{tabular}}\par
\raggedcolumns
\papersec{1}
\partsec{2}{CHEMISTRY}
\sect{SECTION-I (i)}
\qn{1} Which of the following pair of complex both show stereo isomerism as well as structural isomerism.
\begin{qoptions}
\opt{A}{[Pt(en)(CN)$_2$] , [Pt(en)$_3$]Cl$_3$}
\opt{B}{[Co(H$_2$O)$_3$F$_3$] , [Co(NH$_3$)$_4$Cl$_2$]Cl}
\opt{C}{[Pt(NH$_3$)$_2$Cl$_2$Br$_2$] , [Pt(NH$_3$)$_5$Cl)]Br}
\opt{D}{[Co(H$_2$O)$_4$Cl$_2$]Cl.2H$_2$O, [Co(NH$_2$\allowbreak CH$_2$\allowbreak CH$_2$\allowbreak NH$_2$)$_2$\allowbreak Cl$_2$]NO$_3$}
\end{qoptions}
\qn{2} Which of the following is not disproportionation reaction.
\begin{qoptions}
\opt{A}{P$_4$ $\xrightarrow{\text{NaOH}}$}
\opt{B}{HNO$_2$ $\xrightarrow{\Delta}$}
\opt{C}{XeF$_6$ $\xrightarrow{\text{H}_2\text{O/OH}^-\text{(excess)}}$}
\opt{D}{NaOCl $\xrightarrow[\Delta]{\text{dil. HCl}}$}
\end{qoptions}
\qn{3} The rate of effusion of two gases `A' and `B' under identical conditions of temperature and pressure are in the ratio of 2 : 1. What is the ratio of rms velocity of their molecules if temperature of A \& B are in ratio 2 : 1
\begin{qoptions}
\opt{A}{$\sqrt2$ : 1}
\opt{B}{2 : 1}
\opt{C}{1 : $\sqrt2$}
\opt{D}{2$\sqrt2$ : 1}
\end{qoptions}
\qn{4} A(Hydrocarbon) (C $\Rightarrow$ 88.24\%) [Molecular weight of A = 68]\par
A $\xrightarrow[\text{n - propyl bromide}]{\text{Na}}$ B (C$_8$H$_{14}$)\par
B $\xrightarrow[\text{dil.H}_2\text{SO}_4]{\text{Hg}^{2+}}$ C$_5$H$_6$O (C)\par
B $\xrightarrow[\text{H}^+]{\text{KMnO}_4}$ 2-isomeric acids D and E.\par
`A' is :
\begin{qoptions}
\opt{A}{H$_3$C$-$H$_2$C$-$C$\equiv$C$-$CH$_3$}
\opt{B}{CH$_3-$CH$_2-$CH$_2-$C$\equiv$C$-$H}
\opt{C}{CH$_3-$CH$(\mathrm{CH_3})-$C$\equiv$CH}
\opt{D}{\qfigi{p1_chem_s1_q04_D}}
\end{qoptions}
\qn{5} \qfigi{p1_chem_s1_q05_R}\par
(C) is
\begin{qoptions}
\opt{A}{\qfigi{p1_chem_s1_q05_A}}
\opt{B}{\qfigi{p1_chem_s1_q05_B}}
\opt{C}{\qfigi{p1_chem_s1_q05_C}}
\opt{D}{\qfigi{p1_chem_s1_q05_D}}
\end{qoptions}
\sect{SECTION-I (ii)}
\qn{6} Which of the following order is/are correct?
\begin{qoptions}
\opt{A}{RbCl $>$ LiCl (Ionic character)}
\opt{B}{AlCl$_3$ $>$ BCl$_3$ (Lewis acidic strength)}
\opt{C}{Atomic nitrogen $>$ molecular nitrogen (Ionization energy)}
\opt{D}{Li$_2$CO$_3$ $>$ Cs$_2$CO$_3$ (Thermal stability)}
\end{qoptions}
\qn{7} In the homolytic cleavage of N$_2$O$_4$ which of the following propertie(s) changes.
\begin{qoptions}
\opt{A}{Polarity}
\opt{B}{Bond length of NO}
\opt{C}{Bond angle (O $-\,\widehat{\mathrm N}-$ O)}
\opt{D}{Magnetic behaviour}
\end{qoptions}
\qn{8} Which of the following is/are correct at 298K :
\begin{qoptions}
\opt{A}{$\Delta_\mathrm fG^\circ$ (H$_{2(\mathrm g)}$) = 0}
\opt{B}{$\Delta_\mathrm fG^\circ$ (H$_2$O, $\ell$) = 0}
\opt{C}{$\Delta_\mathrm fG^\circ$ (D$_2$, g) = 0}
\opt{D}{$\Delta_\mathrm fG^\circ$ (H$^+$, aq.) = 0}
\end{qoptions}
\qn{9} \qfig{p1_chem_s2_q09_1}
P$_1$ \& P$_2$ will be distinguished by
\begin{qoptions}
\opt{A}{2, 4, D.N.P}
\opt{B}{alkaline Cu$^{2+}$ (tartarate)}
\opt{C}{NaHSO$_3$}
\opt{D}{Schiff's reagent}
\end{qoptions}
\sect{SECTION-II}
\qn{1} Find the sum of element which does not have f-subshell e$^-$.\par
Th , La , Hg , Pb , Os , Pt , Rh
\qn{2} When CrI$_3$ reacts with H$_2$O$_2$ in presence of NaOH forming Na$_2$CrO$_4$, NaIO$_4$ \& H$_2$O then find how many times will be rate of disappearance of NaOH compared to CrI$_3$.
\qn{3} \qfigi{p1_chem_s3_q03_1} $\xrightarrow[\text{I}_2]{\text{NaOH}}$ Number of iodoform molecules produced per molecule.
\qn{4} If pKa values for protonated Lysine molecule are pKa$_1$ = 2.18 ; pKa$_2$ = 8.95 ; pKa$_3$ = 10.79 respectively then , Calculate vlue of (x $-$ 1) where x is iso-electric point of lysine.\par
\qfig{p1_chem_s3_q04_1}
Answer should be nearest integral value
\qn{5} Identify reactions in which new C$-$C bond is formed ?\par
(i) Aldol \quad (ii) Cannizaro \quad (iii) Reimer Tiemann \quad (iv) Hoffmann's Bromamide \quad (v) Sandmeyer reaction \quad (vi) Gattermann Aldehyde sysnthesis \quad (vii) Friedel craft reaction \quad (viii) Wurtz reaction
\qn{6} Find the maximum number of electron(s) in an orbital which having n = 4 and $\ell$ = 2.
\sect{SECTION-IV}
\qn{1} Match the column :\par
{\footnotesize\setlength{\tabcolsep}{1.5pt}\noindent\begin{tabular}{|c|L{3.5cm}|c|L{3.4cm}|}
\hline
 & \multicolumn{1}{>{\centering\arraybackslash}p{3.5cm}|}{Column-I} & & \multicolumn{1}{>{\centering\arraybackslash}p{3.4cm}|}{Column-II} \\ \hline
(A) & Tetragonal and Hexagonal & (P) & are two crystal systems \\ \hline
(B) & Cubic and Rhombohedral & (Q) & a = b $\neq$ c \\ \hline
(C) & Monoclinic and Triclinic & (R) & a $\neq$ b $\neq$ c \\ \hline
(D) & Cubic and Hexagonal & (S) & a = b = c \\ \hline
 & & (T) & $\alpha=\beta=\gamma$ \\ \hline
\end{tabular}}\par
\qn{2} Column-I and Column-II contain data on Schrodinger Wave $-$ Mechanical model, where symbols have their usual meaning. Match the columns :\par
{\footnotesize\setlength{\tabcolsep}{1.5pt}\noindent\begin{tabular}{|c|L{4.3cm}|c|L{1.4cm}|}
\hline
 & \multicolumn{1}{>{\centering\arraybackslash}p{4.3cm}|}{Column-I} & & \multicolumn{1}{>{\centering\arraybackslash}p{1.4cm}|}{Column-II} \\ \hline
(A) & \qfigi{p1_chem_s4_q02_A} & (P) & 4s \\ \hline
(B) & \qfigi{p1_chem_s4_q02_B} & (Q) & 5p$_\mathrm x$ \\ \hline
(C) & $\Psi(\theta,\phi)$ = K (independent of $\theta$ \& $\phi$) & (R) & 3s \\ \hline
(D) & Atleast one angular node is present & (S) & 6d$_\mathrm{xy}$ \\ \hline
 & & (T) & 2p$_\mathrm z$ \\ \hline
\end{tabular}}\par
\raggedcolumns
\papersec{1}
\partsec{3}{MATHEMATICS}
\sect{SECTION-I (i)}
\qn{1} The value of $\displaystyle\lim_{n\to\infty}\sum_{n=1}^{n}$ is equal to -\par
\noindent\resizebox{\linewidth}{!}{$\dfrac{n^2\tan^{-1}\left(\dfrac{2n+1}{1+(n^2+n)^2}\right)+(2n+1)\tan^{-1}(n+1)^2}{(n^2+n)^2\tan^{-1}n^2\cdot\tan^{-1}(n+1)^2}$}\par
\begin{qoptions}
\opt{A}{$-\dfrac\pi4$}
\opt{B}{$\dfrac4\pi$}
\opt{C}{$\dfrac\pi2$}
\opt{D}{$\dfrac2\pi$}
\end{qoptions}
\qn{2} Let A = [a$_{ij}$]$_3$ be an orthogonal matrix, B and C are two other matrices of same order as A of which B is symmetric and C is skew symmetric such that A(B + C) = KI where K is non zero scalar. If Tr.(B) = 8 Tr.(A) $\neq$ 0 then K is -\par
[Note : Tr.(P) denotes trace of matrix P]
\begin{qoptions}
\opt{A}{2}
\opt{B}{3}
\opt{C}{4}
\opt{D}{8}
\end{qoptions}
\qn{3} If one number is selected from each of the three sets A = \{1,2,3,4,5,6,7\}, B = \{1,2,3,4\} and C = \{2,3,4,5,6,7\} then probability that product of the number selected from one and sum of the numbers selected from the other two, is a prime, is -
\begin{qoptions}
\opt{A}{$\dfrac17$}
\opt{B}{$\dfrac47$}
\opt{C}{$\dfrac{13}{56}$}
\opt{D}{$\dfrac1{14}$}
\end{qoptions}
\qn{4} Let $\vec{\mathrm v}_1$ = a$^2\hat{\mathrm i}$ $-$ 2$\hat{\mathrm j}$ + (a + 1)$\hat{\mathrm k}$ and $\vec{\mathrm v}_2$ = 2$\hat{\mathrm i}$ + $\hat{\mathrm j}$ $-$ (a$^2$ $-$ 1)$\hat{\mathrm k}$ be two vectors such that $\vec{\mathrm v}_2$ lies in the plane 3x $-$ $\lambda$y + z = 1 and ($\vec{\mathrm v}_1$ + $\vec{\mathrm v}_2$).$\vec{\mathrm v}_1$ = a$^4$ + a$^2$ + 2a + 5 then greatest volume of the parallelopiped whose co-terminous edges are represented by $\vec{\mathrm v}_1$, $\vec{\mathrm v}_2$ and $\vec{\mathrm n}$ where $\vec{\mathrm n}$ is the normal vector of the given plane, is -
\begin{qoptions}
\opt{A}{3}
\opt{B}{5}
\opt{C}{15}
\opt{D}{25}
\end{qoptions}
\qn{5} Let A = $\left\{\dfrac{x^2-p^2}{x^2-4x+p^2}\,\middle|\,x\in\mathrm R,\ p\in\mathrm R\right\}$ and\par
B = $\left\{p\in\mathrm I\,\middle|\,\dfrac{x^2-p^2}{x^2-4x+p^2}\in\mathrm R\right\}$\par
If A is ($-\infty$, $\infty$) then the value of $|$B$|$ is -\par
[Note : $|$M$|$ denotes the number of elements in set M]
\begin{qoptions}
\opt{A}{2}
\opt{B}{3}
\opt{C}{4}
\opt{D}{5}
\end{qoptions}
\sect{SECTION-I (ii)}
\qn{6} Let y = f(x) satisfies the differential equation 3x$^3$ydx $-$ x$^4$dy = (y$^2$ + x$^3$y)dx where x $>$ 0, y $>$ 0 and x $\neq$ $\dfrac12$. If y(1) = 1 then identify which of the following statement(s) is(are) correct ?
\begin{qoptions}
\opt{A}{Area of the triangle formed by tangent and normal drawn to the curve at (1, 1) with the x-axis 1 sq. unit.}
\opt{B}{If f(x) = 1 + m(x $-$ 1) has three distinct solutions then m $>$ 1.}
\opt{C}{If f(x) = 1 + m(x $-$ 1) has two distinct solutions then m $\in$ [0, 1).}
\opt{D}{If f(x) = 1 + m(x $-$ 1) has two distinct solutions then m $\in$ ($-\infty$, 1).}
\end{qoptions}
\qn{7} Consider H : $\dfrac{x^2}4-\dfrac{y^2}{16}$ = 1 and P : y$^2$ = 4x,\par
If M(x, y) lies inside the parabola from where two distinct tangents are drawn to the hyperbola H, then identify the correct statement(s) -
\begin{qoptions}
\opt{A}{If tangents touches two different branches of the hyperbola then area of the region in which M(x, y) lies is 2/3 sq. units.}
\opt{B}{If tangents touches two different branches of the hyperbola then area of the region in which M(x, y) lies is 1/3 sq. units.}
\opt{C}{If tangents touches two same branches of the hyperbola then area of the region in which M(x, y) (0 $\le$ x $\le$ 2) lies is $\left(\dfrac{16\sqrt2-2}3\right)$ sq. units.}
\opt{D}{If tangents touches two same branches of the hyperbola then area of the region in which M(x, y) (0 $\le$ x $\le$ 2) lies is $\left(\dfrac{8\sqrt2-1}3\right)$ sq. units.}
\end{qoptions}
\qn{8} If f(x) = max(x$^2$ $-$4, 2 $-$ $|$x$|$) and g(x) = f(f(x)) then identify which of the following statement(s) is(are) correct ?
\begin{qoptions}
\opt{A}{g(x) is non-derivable at 5 points.}
\opt{B}{$\displaystyle\int_{-1}^{2}$ g(x)dx = $\dfrac12$}
\opt{C}{The number of solution of the equation g(x) = $\dfrac{2015}{2014}$ is 8.}
\opt{D}{g(x) is continuous for all x $\in$ R.}
\end{qoptions}
\qn{9} If maximum length of the tangent drawn to the circle which touches the straight line x + y = 4 at (2,2) and also the positive co-ordinate axes, from the point (x, $\pm$ (4 $-$ $|$x$|$)) where $-$ 6 $\le$ x $\le$ 4 is 4$\sqrt{\mathrm a+\mathrm b\sqrt2}$ then (a + b) is\par
(Where a, b $\in$ Z)
\begin{qoptions}
\opt{A}{11}
\opt{B}{13}
\opt{C}{7}
\opt{D}{15}
\end{qoptions}
\sect{SECTION-II}
\qn{1} If\par
\noindent\resizebox{\linewidth}{!}{$\displaystyle\int_{\pi/6}^{7\pi/6}\left(\ell\mathrm n\left|3\cos\left(x+\dfrac\pi3\right)+5\cos x\right|+\left|3\sin^2x-6\sin x\cos x+\cos^2x-2\right|\right)dx=\pi\ell\mathrm np+\sqrt q$}\par
where p, q $\in$ Q then the value of (p$\cdot$q) is equal to -
\qn{2} Let f : R$\to$R be a function such that f(x$^5$ + x $-$ 2) = x$^3$ $-$ 1 $\forall$ x $\in$ R and g be the inverse function of f.\par
If $\ell$ = $\left.\dfrac{d}{dx}\left((x^3+1)\cdot g(f^3(g(f(g(x)))))\right)\right|_{x=7^{1/3}}$ and m = $\left.\dfrac{d^2}{dx^2}\left(\dfrac{f(f(g(g(f(x)))))}{x^2}\right)\right|_{x=-2}$ then the value of ($\ell$ (7$^{-2/3}$) + 8m) is -
\qn{3} If $\alpha$ is a non-real 7$^{\mathrm{th}}$ root of 3 and x = $\alpha$ + $\alpha^2$ then the value of x$^7$ $-$ 21[x$^2$(x + 1) + x$^2$ + x + 1] + 15 is -
\qn{4} Consider f(x) = $||$x $-$ 2$|$ $-$ $|$3x $-$ 1$|$ + $|$2x + 1$||$ and\par
g(x) = $\left[\left|-x^2+\dfrac32x+1\right|\right]$\par
[Note : [y] denotes greatest integer less than or equal to y.]\par
If m and n are the no. of points where h(x) = f(x)$\cdot$g(x) is discontinuous and non-derivable respectively then (m + n) is -
\qn{5} If x + y + z + w = a$^2$ + b$^2$ $-$ 4a + 4b where a$^2$ + b$^2$ = 4 and a, b $\in$ R then number of quadruplets (x, y, z, w) where x,y,z,w $\in$ N is -
\qn{6} Let the locus of the centre of the family of circles which intersects the parabola x$^2$ $-$ 2y + 3x = 4 orthogonally at the point (1/2 , $-$9/8) be C and $\ell$ be the least distance of the curve C from y = x $-$ 1/x (x $<$ 0) then (8$\sqrt5\,\ell$) is
\sect{SECTION-IV}
\qn{1} {\footnotesize\setlength{\tabcolsep}{1.5pt}\noindent\begin{tabular}{|c|L{3.7cm}|c|L{1.5cm}|}
\hline
 & \multicolumn{1}{>{\centering\arraybackslash}p{3.7cm}|}{Column-I} & & \multicolumn{1}{>{\centering\arraybackslash}p{1.5cm}|}{Column-II} \\ \hline
(A) & If the point ($\alpha$, $\alpha^2-3$), $\alpha\in$ R lies in acute angular region between the lines 7x $-$ y $-$ 15 = 0 and x $-$ y + 3 = 0 then largest integral value of $\alpha$ is & (P) & 2 \\ \hline
(B) & If A, B, C are three sets such that n(A$\cap$B) = n(B$\cap$C) = n(C$\cap$A) = n(A$\cap$B$\cap$C) = 2 then n((A$\times$B)$\cap$(B$\times$C)) is & (Q) & 3 \\ \hline
(C) & If the shortest distance between the lines formed by the intersection of the plane 2x $-$ 3y + z = 1 with the planes x $-$ y $-$ z = 1 and x $-$ y $-$ z + 1 = 0 is 2$\sqrt{\dfrac ab}$ where a and b are prime numbers then $|$b $-$ a$|$ is & (R) & 4 \\ \hline
(D) & Let 3 male and 4 female teachers be ready for a ramp walk in a line, if the number of ways in which even number of female teachers are there after every male teacher in the line, is 60 (K!) then K is & (S) & 5 \\ \hline
 & & (T) & 6 \\ \hline
\end{tabular}}\par
\qn{2} {\footnotesize\setlength{\tabcolsep}{1.5pt}\noindent\begin{tabular}{|c|L{3.7cm}|c|L{1.5cm}|}
\hline
 & \multicolumn{1}{>{\centering\arraybackslash}p{3.7cm}|}{Column-I} & & \multicolumn{1}{>{\centering\arraybackslash}p{1.5cm}|}{Column-II} \\ \hline
(A) & Let $\alpha$ and $\beta$ be two positive integral roots of the equation x$^3$ + (221 $-$ 4m)x$^2$ + 220x + 4m = 0 where $\alpha$ $<$ $\beta$ $<$ 20. Then $\left(\dfrac{\alpha+\beta+\mathrm m}{19}\right)$ is less than or equal to - & (P) & 0 \\ \hline
(B) & If the set of values of p for which the equation sin$^{-1}$(2$|$x$|$ $-$ 1) = p($|$x$|$ $-$ 1/2) has 6 distinct solutions is (a, b) then (a + b) is less than or equal to - & (Q) & 2 \\ \hline
(C) & Let $\bar{\mathrm x}_\omega$ be the weighted mean of first n natural numbers when their weights are equal to their squares respectively then $\displaystyle\lim_{n\to\infty}\dfrac{\bar{\mathrm x}_\omega}{\mathrm n}$ is less than or equal to & (R) & 3 \\ \hline
(D) & If f is polynomial such that f(f(f(x))) $-$ 2f(f(x)) + 3x $-$ 7 = 0 $\forall$ x $\in$ R then $\dfrac12\displaystyle\int_0^1$ (f(x) $-$ f$^{-1}$(x))dx is less than or equal to - & (S) & 5 \\ \hline
 & & (T) & 8 \\ \hline
\end{tabular}}\par
\raggedcolumns
\papersec{2}
\partsec{1}{PHYSICS}
\sect{SECTION-I (i)}
\qn{1} Figure shows a circuit consisting of an inductor and resistors connected with a cell. All ammeter are ideal. The switch was closed at t = 0. Choose the correct statement(s).\par
\qfig{p2_phy_s1_q01_1}
\begin{qoptions}
\opt{A}{Reading of A$_1$ is sum of reading of A$_2$ and A$_3$.}
\opt{B}{Reading of A$_2$ will be $\dfrac{\varepsilon}{10\mathrm R}\left(4+\mathrm e^{-\frac{5\mathrm{Rt}}{2\mathrm L}}\right)$.}
\opt{C}{Reading of A$_3$ will be $\dfrac{\varepsilon}{5\mathrm R}\left(1-\mathrm e^{-\frac{5\mathrm{Rt}}{2\mathrm L}}\right)$.}
\opt{D}{In steady state reading of A$_1$ will be $\dfrac{\varepsilon}{2\mathrm R}$.}
\end{qoptions}
\qn{2} The ground state and first excited state energies of hydrogen atom are $-$13.6 eV and $-$3.4 eV respectively. If potential energy in ground state is taken to be zero, then :-
\begin{qoptions}
\opt{A}{Potential energy in the first excited state would be 20.4 eV.}
\opt{B}{Total energy in the first excited state would be 23.8 eV.}
\opt{C}{Kinetic energy in the first excited state would be 3.4 eV.}
\opt{D}{Total energy in the ground state would be 13.6 eV.}
\end{qoptions}
\qn{3} A particle of mass m, charge q is projected from origin with velocity v$_0\hat{\mathrm i}$ in a region where uniform magnetic field $\vec{\mathrm B}$ = B$_0$($-\hat{\mathrm k}$) exist. If it experience drag force equal to $\vec{\mathrm F}$ = $-$K$\vec{\mathrm v}$, where K is a positive constant and $\vec{\mathrm v}$ is velocity vector, then :-\par
\qfig{p2_phy_s1_q03_1}
\begin{qoptions}
\opt{A}{Distance travelled by particle before its velocity is along the positive x-axis is $\dfrac{\mathrm{mv}_0}{\mathrm K}\left(1-\mathrm e^{-\frac{2\pi\mathrm K}{\mathrm{qB}}}\right)$.}
\opt{B}{Distance travelled by particle before its velocity is along the positive x-axis is $\dfrac{2\mathrm{mv}_0}{\mathrm K}\left(1-\mathrm e^{-\frac{2\pi\mathrm K}{\mathrm{qB}}}\right)$.}
\opt{C}{Speed of the particle decreases exponentially with time.}
\opt{D}{Particle finally stops at origin.}
\end{qoptions}
\qn{4} Consider two solid spherical asteroid of uniform density of mass M and radius R. In one asteroid a tunnel of very small size of depth R is bored to the centre and in other asteroid a spherical cavity of radius $\dfrac{\mathrm R}2$ is made as shown in the figure. Now, identical particles of mass m are dropped into cavities of both asteroids from top most point P. The force experienced by particles are F$_\mathrm I$ and F$_\mathrm{II}$ respectively in cavities of asteroid I and II, when they are x distance away from centre of asteroids. If the time taken by particles to reach the centre of asteroids is T$_\mathrm I$ and T$_\mathrm{II}$ respectively then\par
\qfig{p2_phy_s1_q04_1}
\begin{qoptions}
\opt{A}{The ratio of F$_\mathrm I$/F$_\mathrm{II}$ is equal to 2x/R.}
\opt{B}{The ratio of F$_\mathrm I$/F$_\mathrm{II}$ is equal to x/R.}
\opt{C}{The ratio of T$_\mathrm I$/T$_\mathrm{II}$ is equal to $\pi$/4.}
\opt{D}{The ratio of T$_\mathrm I$/T$_\mathrm{II}$ is equal to $\pi$/2.}
\end{qoptions}
\qn{5} Four identical conducting plates of area A each are connected as shown in circuit and the gap `d' between adjacent plates is same. Initially charge on each plate is zero. Select the correct statements.\par
\qfig{p2_phy_s1_q05_1}
\begin{qoptions}
\opt{A}{Equivalent capacitance between the terminals of battery is $\left(\dfrac{3\epsilon_0\mathrm A}{2\mathrm d}\right)$.}
\opt{B}{Sum of charges on plate-1 and plate-4 is zero.}
\opt{C}{The charge on right side of plate-2 is $\left(-\dfrac{\epsilon_0\mathrm{AV}}{\mathrm d}\right)$.}
\opt{D}{The charge on left side of plate-2 is zero.}
\end{qoptions}
\qn{6} Train is moving with constant speed along a circular track. If length of the train is one fourth of length of circular track then which of the following is/are correct options? (Assume that sound source is at engine).
\begin{qoptions}
\opt{A}{Frequency observed by a passenger who is sitting in the middle of train (equidistance from front and rear end) will continuously increase.}
\opt{B}{Frequency observed by a passenger who is sitting in the middle of train (equidistance from front and rear end) will remain constant but more than actual frequency.}
\opt{C}{Frequency observed by passenger who is sitting in the middle of train (equidistance from front and rear end) will remains constant and equal to actual frequency.}
\opt{D}{Wavelength observed by the person who is on the rear end of train is more than the actual wavelength of sound wave.}
\end{qoptions}
\qn{7} A particle of mass m moving in XY horizontal plane strikes the end of the vertical rod at an angle 37$^\circ$ with Y-axis. Rod has mass m and length $\ell$ and it is free to rotate about hinge along Y-axis. If particle strikes rod with speed u and sticks to it, find the CORRECT choice(s).\par
\qfig{p2_phy_s1_q07_1}
\begin{qoptions}
\opt{A}{Velocity of particle just after impact will be $-\dfrac{4\mathrm u\hat{\mathrm i}}5+\dfrac{9\mathrm u\hat{\mathrm j}}5$.}
\opt{B}{Velocity of particle just after impact will be $-\dfrac{9\mathrm u\hat{\mathrm i}}{20}+\dfrac{4\mathrm u\hat{\mathrm j}}{10}$.}
\opt{C}{Minimum speed so that rod completes full rotation about hinge is $\dfrac{10}3\sqrt{2\mathrm g\ell}$.}
\opt{D}{Angular momentum of system during collision is conserved about Y-axis.}
\end{qoptions}
\sect{SECTION-I (ii)}
\paraq{8}{9}
A large, flat sheet carries a uniformly distributed electric current with current per unit width J. This current creates a magnetic field on both sides of the sheet, parallel to the sheet and perpendicular to the current, with magnitude B = $\dfrac12\mu_0\mathrm J_\mathrm s$. If the current is in the y direction and oscillates in time according to J$_{\max}$ (cos $\omega$t)$\hat{\mathrm j}$, the sheet radiates an electromagnetic wave. Figure shows such a wave emitted from one point on the sheet chosen to be the origin. Such electromagnetic waves are emitted from all points on the sheet. The magnetic field of the wave to the right of the sheet is described by the wave function\par
$\vec{\mathrm B}$ = $\dfrac12\mu_0\mathrm J_{\max}$ [cos(kx $-$ $\omega$t)]$\hat{\mathrm k}$\par
\qfig{p2_phy_s2_para8_9_1}
\qn{8} The electric field of wave to the right of the sheet is :-
\begin{qoptions}
\opt{A}{$\dfrac12\mu_0\mathrm{cJ}_{\max}\cos(\mathrm{kx}-\omega\mathrm t)\hat{\mathrm j}$}
\opt{B}{$-\dfrac12\mu_0\mathrm{cJ}_{\max}\cos(\mathrm{kx}-\omega\mathrm t)\hat{\mathrm j}$}
\opt{C}{$\dfrac12\dfrac{\mu_0\mathrm J_{\max}}{\mathrm c}\cos(\mathrm{kx}-\omega\mathrm t)\hat{\mathrm k}$}
\opt{D}{$-\dfrac12\dfrac{\mu_0\mathrm J_{\max}}{\mathrm c}\cos(\mathrm{kx}-\omega\mathrm t)\hat{\mathrm k}$}
\end{qoptions}
\qn{9} If the intensity of radiation emitted by the sheet in each direction is 570 W/m$^2$, the maximum value of sinusoidal current density is :-
\begin{qoptions}
\opt{A}{1.28 A/m}
\opt{B}{2.46 A/m}
\opt{C}{3.48 A/m}
\opt{D}{4.92 A/m}
\end{qoptions}
\paraq{10}{11}
3.2g of oxygen starts from the initial state 1 (p$_1$ = 1.0 atm, V$_1$ = 1.0L), reaches state 2 (V$_2$ = 2.0L) through an isobaric process, and then passes through the isochoric process. The process reaches state 3 (p$_3$ = 2.0atm), and then reaches state 4 through an adiabatic process. The temperature of state 4 is exactly the same as the temperature of state 1. Finally, it returns to state 1 through an isothermal process, forming a cyclic process.\par
\qfig{p2_phy_s2_para10_11_1}
\qn{10} The efficiency of cyclic process is :-
\begin{qoptions}
\opt{A}{1 $-$ $\dfrac{\ell\mathrm n2}2$}
\opt{B}{1 $-$ $\dfrac{3\ell\mathrm n2}9$}
\opt{C}{1 $-$ $\dfrac{6\ell\mathrm n2}{13}$}
\opt{D}{1 $-$ $\dfrac{12\ell\mathrm n2}{17}$}
\end{qoptions}
\qn{11} The variation of entropy with temperature is given by :-
\begin{qoptions}
\opt{A}{\qfigi{p2_phy_s2_q11_A}}
\opt{B}{\qfigi{p2_phy_s2_q11_B}}
\opt{C}{\qfigi{p2_phy_s2_q11_C}}
\opt{D}{\qfigi{p2_phy_s2_q11_D}}
\end{qoptions}
\sect{SECTION-II}
\qn{1} Three identical uniform dominos are arranged as shown in the diagram. The possible values of x such that the dominoes are in stable equilibrium is $\alpha$ $<$ x $<$ $\beta$. $\alpha$, $\beta$ and x are measured in mm. Then the value of $\alpha$ + $\beta$ (in mm) is :-\par
\qfig{p2_phy_s3_q01_1}
\qn{2} A magnetic disc is moving on an air-cushioned table above a metal sheet. Due to the generated eddy currents there is a retarding force exerted on the disc, which is proportional to the speed of the disc. Moving above a sheet made of aluminum the disc stops after covering a distance of 30 cm, but when the sheet is made of copper the disc stops after covering only 20 cm. How much total distance (in cm) will the disc cover if first it travels above a piece of copper sheet of width 15 cm, and then continues its motion above an aluminum sheet? (The initial speeds of the disc are the same in all three cases.)
\qn{3} In the young"s double slit experiment a point source of $\lambda$ = 5000 {\AA} is placed slightly off the central axis as shown in the figure. If 70$^{\mathrm{th}}$ order maximum is formed at P, the value of n is :-\par
\qfig{p2_phy_s3_q03_1}
\qn{4} A rocket is launched normal to the surface of the Earth, away from the Sun along the line joining the Sun and the Earth. The Sun is 3 $\times$ 10$^5$ times heavier than the Earth and is at a distance 2.5 $\times$ 10$^4$ times larger than the radius of the Earth. The escape velocity from Earth's gravitational field is v$_\mathrm e$ = 11.2 kms$^{-1}$. The minimum initial velocity (v$_\mathrm s$) (in km/sec) required for the rocket to be able to leave the Sun-Earth system is closest to :- (Ignore the rotation and revolution of the Earth and the presence of any other planet)
\qn{5} Troposphere is the lowest layer of atmosphere, whose temperature decreases linearly with height (y) according to T = T$_0$ $-$ $\alpha$y, where T$_0$ is temperature at Earth's surface and the constant $\alpha$ is called lapse rate of temperature. The average molar mass of normal atmospheric air is M$_0$ kg/mole and for the height of troposphere, the gravity field g is constant. The atmospheric pressure at ground is P$_0$ and gas constant is R.\par
The atmospheric pressure as a function of height y will be P (y) = P$_0\left(1-\dfrac{\alpha\mathrm y}{\mathrm T_0}\right)^{\frac{\mathrm{nM}_0\mathrm g}{2\alpha\mathrm R}}$, then value of n is :-
\qn{6} The web of negligible weight has the shape shown in the picture. The hexagons are regular and divide the radial fibers into equal parts of length $\ell$. The ends of the web are fixed at the same height above the ground. At the beginning, the web is not sagged. When a spider of mass m moves to the center of the web, the web bends so that its center moves closer to the ground by a distance h. Determine the Young's modulus of elasticity Y of a spider fiber that has a diameter d. If the answer is x $\times$ 10$^{17}$ Pa, fill the value of x. (Given : $\ell$ = 0.10 m, m = 0.010 kg, h = 5.0 mm, d = 0.15 $\mu$m)\par
\qfig{p2_phy_s3_q06_1}
\raggedcolumns
\papersec{2}
\partsec{2}{CHEMISTRY}
\sect{SECTION-I (i)}
\qn{1} The chain of reaction is completely hypothetical and based upon formula matching only.\par
Consider 1 molecule of all the compounds H$_2$O, SO$_3$, SO$_2$ , A, B, C, D, E in the given chain reaction\par
\qfig{p2_chem_s1_q01_1}
Select the CORRECT statement :
\begin{qoptions}
\opt{A}{When one oxygen atom is removed from formula of B then formula of C is formed}
\opt{B}{C has S $-$ S linkage in it's structure}
\opt{C}{E has S $-$ O $-$ S linkage in it's structure.}
\opt{D}{C is H$_2$S$_2$O$_6$ and E is H$_2$S$_2$O$_5$}
\end{qoptions}
\qn{2} {\footnotesize\setlength{\tabcolsep}{1.5pt}\noindent\begin{tabular}{L{3.4cm}L{3.6cm}}
Co-ordination entity & Wave length of light absorbed (nm) \\
{}[CoCl(NH$_3$)$_5$]$^{2+}$ & P \\
{}[Co(NH$_3$)$_5$(H$_2$O)]$^{3+}$ & Q \\
{}[Co(NH$_3$)$_6$]$^{3+}$ & R \\
{}[Co(CN)$_6$]$^{3-}$ & S \\
{}[Cu(H$_2$O)$_4$]$^{2+}$ & T \\
\end{tabular}}\par
Select the correct orders.
\begin{qoptions}
\opt{A}{P $>$ R $>$ Q}
\opt{B}{T $>$ P $>$ Q $>$ R $>$ S}
\opt{C}{P $>$ Q $>$ R}
\opt{D}{T $>$ Q $>$ P $>$ R $>$ S}
\end{qoptions}
\qn{3} Which reactions/s can produce 2$^\circ$-amines as major product?
\begin{qoptions}
\opt{A}{\qfigi{p2_chem_s1_q03_A}}
\opt{B}{\qfigi{p2_chem_s1_q03_B}}
\opt{C}{\qfigi{p2_chem_s1_q03_C}}
\opt{D}{\qfigi{p2_chem_s1_q03_D}}
\end{qoptions}
\qn{4} \qfigi{p2_chem_s1_q04_1}\par
Correct statement(s) for above reaction sequence :
\begin{qoptions}
\opt{A}{Product (P) shows Tollen's test.}
\opt{B}{Product (P) shows Iodoform test.}
\opt{C}{Above reaction sequence involve elimination reaction.}
\opt{D}{Product (P) can show Lucas test}
\end{qoptions}
\qn{5} Which of the following reaction/s is/are correctly represented with major products?
\begin{qoptions}
\opt{A}{\qfigi{p2_chem_s1_q05_A}}
\opt{B}{\qfigi{p2_chem_s1_q05_B}}
\opt{C}{\qfigi{p2_chem_s1_q05_C}}
\opt{D}{\qfigi{p2_chem_s1_q05_D}}
\end{qoptions}
\qn{6} P$_\mathrm A$ = (235 y $-$ 125 xy) mm of Hg. P$_\mathrm A$ is partial pressure of A, x is mole fraction of B in liquid phase in the mixture of two liquids A and B and y is the mole fraction of A in vapour phase, then :
\begin{qoptions}
\opt{A}{P$^\circ_\mathrm A$ = 110}
\opt{B}{P$^\circ_\mathrm B$ = 110}
\opt{C}{P$^\circ_\mathrm A$ = 235}
\opt{D}{P$^\circ_\mathrm B$ = 235}
\end{qoptions}
\qn{7} The correct information(s) is/are :
\begin{qoptions}
\opt{A}{The minimum concentration of an electrolyte in millimoles per litre required to cause precipitation of a sol in two hours is called coagulation power}
\opt{B}{Smaller the coagulation value, higher will be the coagulation power of that ion.}
\opt{C}{Lyophilic sols may be coagulated by adding an electrolyte or a suitable solvent.}
\opt{D}{Lyophobic colloids used to protect lyophilic colloids are called protective colloids.}
\end{qoptions}
\sect{SECTION-I (ii)}
\paraq{8}{9}
The electrochemical cell shown below is a concentration cell.\par
M $|$ M$^{2+}$ (saturated solution of a sparingly soluble salt, MX$_2$) $||$ M$^{2+}$ (0.001 mol dm$^{-3}$) $|$ M\par
The emf of the cell depends on the difference in concentrations of M$^{2+}$ ions at the two electrodes. The emf of the cell at 298 K is 0.059V.
\qn{8} The value of $\Delta$G (kJ mol$^{-1}$) for the given cell M(s) + M$^{2+}$(aq) $\longrightarrow$ M(s) + M$^{2+}$ (aq) is (take F = 96500 C mol$^{-1}$)
\begin{qoptions}
\opt{A}{$-$5.7}
\opt{B}{5.7}
\opt{C}{11.4}
\opt{D}{$-$11.4}
\end{qoptions}
\qn{9} The solubility product (K$_\mathrm{sp}$ ; mol$^3$ dm$^{-9}$) of MX$_2$ at 298 K based on the information available for the given concentration cell is x $\times$ 10$^{-15}$ (take 2.303 $\times$ R $\times$ 298$|$F = 0.059 V). The value of x is ?
\begin{qoptions}
\opt{A}{1 $\times$ 10$^{-15}$}
\opt{B}{4 $\times$ 10$^{-15}$}
\opt{C}{1 $\times$ 10$^{-12}$}
\opt{D}{4 $\times$ 10$^{-12}$}
\end{qoptions}
\paraq{10}{11}
If values of quantum numbers can be considered as follows.\par
n = 1, 2, 3 ....... $\infty$\par
$\ell$ = 0 to (n + 1)\par
m = $-\ell$ to +$\ell$ (integral value including zero)\par
s = $-\dfrac14$, $-\dfrac12$, +$\dfrac12$, +$\dfrac14$
\qn{10} According to above consideration, the maximum number of e$^-$ associated to 2$^{\mathrm{nd}}$ shell is :
\begin{qoptions}
\opt{A}{64}
\opt{B}{32}
\opt{C}{16}
\opt{D}{8}
\end{qoptions}
\qn{11} According to above consideration, the maximum number of e$^-$ associated to a p-orbital is :
\begin{qoptions}
\opt{A}{2}
\opt{B}{4}
\opt{C}{6}
\opt{D}{12}
\end{qoptions}
\sect{SECTION-II}
\qn{1} How many of the following cations gives ppt with H$_2$S in acidic solution :-\par
Zn$^{2+}$ , Cu$^{2+}$ , Ni$^{2+}$ , Sr$^{2+}$ , Al$^{3+}$ , Bi$^{3+}$ , Sn$^{2+}$ , Pb$^{2+}$
\qn{2} Among the following statements\par
(i) Ag$_2$SO$_3$ is insoluble in dilute HNO$_3$.\par
(ii) Nitrogen dioxide (NO$_2$) is absorbed by ferrous sulphate solution forming brown solution.\par
(iii) Heavy metal chlorides like AgCl, HgCl$_2$. etc. do not respond to chromyl chloride test.\par
(iv) The group reagent for III group basic radicals is NaOH.\par
(v) Green colour of alkaline tetrahydroxidochromate (III) becomes yellow on adding H$_2$O$_2$.\par
(vi) Titan yellow solution gives blue colouration/ppt. with a acidic solution containing Mg$^{2+}$ ions\par
(vii) Solution of nitrite is decomposed by sulphamic acid.\par
(viii) Fe$^{2+}$ ions give brown colour precipitate with [Fe(CN)$_6$]$^{3-}$ ions solution.\par
(ix) Green precipitate of Cr(OH)$_3$ is soluble in excess Na$_2$O$_2$.\par
If ``a'' is total number of correct while ``b'' is total number of incorrect statements find $\dfrac{\mathrm a^2+\mathrm b^2}3$ .
\qn{3} \qfig{p2_chem_s3_q03_S}
Number of compounds that gives Tollen's test = x.\par
Number of compounds that gives Fehling's test = y.\par
Then x + y = \underline{\hspace{1.5cm}}
\qn{4} `X' is the number of possible products (including stereoisomers) in the following reaction, then value of $\dfrac{\mathrm X}2$ is\par
\qfig{p2_chem_s3_q04_1}
\qn{5} Fixed amount of an ideal gas contained in a sealed regid vessel (V = 24.6 litre) at 1. 0 bar is heated reversibly from 27$^\circ$C to 127$^\circ$ C. Determine change in Gibb's energy (in J), if entropy of gas S = 10 + 10$^{-2}$ T (J/K)
\qn{6} The solubility of a salt (AB) of weak acid at pH = 3 is $\sqrt{\mathrm Y}\times10^{-3}$ mol L$^{-1}$. The value of Y is\underline{\hspace{0.8cm}}.\par
(Given that the value of solubility product of AB (K$_\mathrm{sp}$) = 2 $\times$ 10$^{-10}$ and the value of ionization constant of HB(K$_\mathrm a$) = 1 $\times$ 10$^{-8}$)
\raggedcolumns
\papersec{2}
\partsec{3}{MATHEMATICS}
\sect{SECTION-I (i)}
\qn{1} Let A, B, C be 3 sets of complex numbers as defined below :\par
A = \{z : $|$z + 1$|$ $\le$ 2 + Re(z)\}, B = \{z : $|$z $-$ 1$|$ $\ge$ 1\} and C = $\left\{z:\left|\dfrac{z-1}{z+1}\right|\ge1\right\}$\par
Then which of the following are correct :
\begin{qoptions}
\opt{A}{Area of the region bounded by A $\cap$ B $\cap$ C is $\sqrt3$}
\opt{B}{Area of the region bounded by A $\cap$ B $\cap$ C is 2$\sqrt3$}
\opt{C}{The real part of the complex number in region A $\cap$ B $\cap$ C having maximum amplitude and modulus is $-\dfrac32$}
\opt{D}{The number of point having integral coordinates in region A $\cap$ B $\cap$ C is 7}
\end{qoptions}
\qn{2} If A is the area enclosed between the curves y = $\ell$n(x + e), x = $\ell$n$\left(\dfrac1y\right)$ and x-axis then :
\begin{qoptions}
\opt{A}{A is prime number}
\opt{B}{A is composite number}
\opt{C}{A is irrational number}
\opt{D}{A$^2$ is even integer}
\end{qoptions}
\qn{3} The solution of the differential equation 3$\dfrac{dx}{dy}$ = $\dfrac{x}{x^3-y}$ is x$^\ell$ = mx$^\mathrm n$y + c, then which of the following options are correct :
\begin{qoptions}
\opt{A}{$\ell$ + m + n = 11}
\opt{B}{$\ell$ + n = 9}
\opt{C}{$\ell$ + 2m = 10}
\opt{D}{m + n = 4}
\end{qoptions}
\qn{4} The value of $\displaystyle\sum_{n=2}^{\infty}\dfrac{n^4+3n^2+10n+10}{2^n(n^4+4)}$ is less than :
\begin{qoptions}
\opt{A}{$\dfrac{11}7$}
\opt{B}{$\dfrac{11}8$}
\opt{C}{$\dfrac{11}9$}
\opt{D}{$\dfrac{11}{10}$}
\end{qoptions}
\qn{5} Tangent drawn at point P(2, $-$5) of a parabola intersects its tangent at vertex at (1, $-$3). If (8, 3) is a point on SP, where S is focus of the parabola then
\begin{qoptions}
\opt{A}{focus of the parabola is (5, $-$1)}
\opt{B}{Directrix of the parabola is x $-$ y $-$ 2 = 0}
\opt{C}{Slope of normal to parabola at other end of focal chord drawn from point P is $-$2}
\opt{D}{Circle drawn on PS as diameter touches the straight line x $-$ y $-$ 4 = 0}
\end{qoptions}
\qn{6} Consider the equation (x$^2$ $-$ 4x + 6)$^2$ $-$ (k + 2)(x$^2$ $-$ 4x + 6) + k + 5 = 0 (k $\in$ R). Which of the following are correct :
\begin{qoptions}
\opt{A}{there is no integral value of k for which given equation has exactly 4 real solutions.}
\opt{B}{there is only one integral value of k for which given equation has exactly 3 real solutions.}
\opt{C}{there is only one integral value of k for which given equation has exactly one real solution.}
\opt{D}{there are only 3 positive integral values of k for which given equation has no real solution.}
\end{qoptions}
\qn{7} Consider the limit $\displaystyle\lim_{x\to0}\dfrac{a\tan x+b\tan^{-1}x+c\sin^{-1}x}{x^5}$ = 3, where a, b, c $\in$ R, then correct option/s is/are
\begin{qoptions}
\opt{A}{a = 30}
\opt{B}{b = 10}
\opt{C}{c = $-$40}
\opt{D}{b = $-$5}
\end{qoptions}
\sect{SECTION-I (ii)}
\paraq{8}{9}
Let f(x) = $\begin{cases}x^2-3x+2 & 0\le x\le3\\ |x-5| & 3<x\le8\\ \log_2(12-x) & 8<x\le11\end{cases}$
\qn{8} If f(x) = k has 5 distinct solutions then sum of integral values of k is
\begin{qoptions}
\opt{A}{1}
\opt{B}{3}
\opt{C}{5}
\opt{D}{7}
\end{qoptions}
\qn{9} Let $\ell$ = maximum number of distinct roots of $|$f(x)$|$ = k and m = $\displaystyle\lim_{x\to3}\left[\mathrm f(3)^{\mathrm f(x)}\right]$, then $\ell$ + m is (Where [.] denotes greatest integer function)
\begin{qoptions}
\opt{A}{11}
\opt{B}{9}
\opt{C}{7}
\opt{D}{5}
\end{qoptions}
\paraq{10}{11}
Let $\dfrac d{dx}$(f(x)) = $\dfrac{\mathrm e^{\sin x}}x$, x $>$ 0, then answer the following :
\qn{10} If $\displaystyle\int_1^4\dfrac{2\mathrm e^{\sin x^2}}x\,dx$ = f(k) $-$ f(1) and xf$''$(x) + f$'$(x) = g(x), then k + $\left[\mathrm g\left(\dfrac\pi2\right)\right]$ is
\begin{qoptions}
\opt{A}{9}
\opt{B}{16}
\opt{C}{27}
\opt{D}{36}
\end{qoptions}
\qn{11} Let I = $\displaystyle\int_0^{\pi/2}$ x cos x f$'$(x) dx = e$^\lambda$ $-$ 1 and xf$'$(x) = T(x), then [$\lambda$] + T$\left(\dfrac\pi2\right)$ is equal to : (Where [.] represents G.I.F.)
\begin{qoptions}
\opt{A}{1 + e}
\opt{B}{3}
\opt{C}{2}
\opt{D}{4}
\end{qoptions}
\sect{SECTION-II}
\qn{1} If k is the number of four digit natural numbers, such that repetition of digits is not allowed and absolute difference between its first and last digit is 2 then $\dfrac k{480}$ is :
\qn{2} If a is the minimum value of x$_1$ and b is the maximum value of x$_2$ such that $\displaystyle\int_{x_1}^{x_2}$ (x$^4$ $-$ 2x$^2$) dx is minimum then the area of region enclosed between the curves y = x$^4$ $-$ 2x$^2$, the x-axis and the ordinates x = a, x = b is :
\qn{3} If a variable circle touches the straight line y = x and cuts the chord of length 2$\sqrt2$ from x + y = 0 then the length of latus rectum of the locus of the centers of the circles is :
\qn{4} If the value of\par
\noindent\resizebox{\linewidth}{!}{$\displaystyle\lim_{n\to\infty}$ n$\left(\dfrac1{3n^2+8n+4}+\dfrac1{3n^2+16n+16}+\dfrac1{3n^2+24n+36}+\cdots\right)$}\par
is $\dfrac14\ell$n$\left(\dfrac nk\right)$ where gcd (n, k) = 1 then $\dfrac nk$ is :
\qn{5} For two data sets each of the size 5, variance are given to be 4 and 5 respectively and mean are 2, 4 respectively, the variance of combined data is equal to
\qn{6} Let A is a 2 $\times$ 2 matrix such that AA$^\mathrm T$ = I, where I is identity matrix then the maximum value of the sum of the absolute values of entries of matrix A is K then find [K] where [.] represents greatest integer function :
\end{multicols}
\end{document}
